\documentclass[10pt,a4paper]{amsart}
\usepackage[margin=19mm]{geometry}
\usepackage{amsmath,amssymb,mathtools}
\usepackage[scr=rsfs,cal=euler]{mathalfa}
\usepackage{xfrac}
\usepackage[unicode,hidelinks,bookmarks=false]{hyperref}
\newcommand{\ie}{i.e.\ }
\newcommand{\cf}{cf.\ }
\newcommand{\resp}{respectively\ }
\newcommand{\Subsection}{\subsection}
\providecommand{\nobreakdash}{\nobreak\hbox{-}}
\setlength{\emergencystretch}{2em}
\multlinegap0pt
\usepackage{bm}
\usepackage{bbm}
\usepackage{dsfont}
\usepackage{xspace}
\usepackage{cancel}
\usepackage{mathtools}
\usepackage{amsthm}
\makeatletter
\@ifundefined{thm}{\newtheorem{thm}{Theorem}}{}
\@ifundefined{lem}{\newtheorem{lem}[thm]{Lemma}}{}
\makeatother
\newcommand{\Z}{\mathbb{Z}}
\newcommand{\dlog}{\operatorname{dlog}}
\newcommand{\e}{\varepsilon}
\newcommand{\z}{\zeta}
\title[Elements in $K\_4$ and regulator maps of Fermat curves]{Elements in $K\_4$ and regulator maps of Fermat curves}
\author{Fran\c{c}ois Brunault, David T.-B. G. Lilienfeldt, Yusuke Nemoto}
\date{Source: arXiv:2606.24532v1 (CC BY 4.0)}
\begin{document}
\maketitle

\noindent\emph{Source and licence.} Elements in $K_4$ and regulator maps of Fermat curves. Version arXiv:2606.24532v1 (CC BY 4.0), distributed under the Creative Commons Attribution 4.0 International licence, as recorded in the arXiv licence metadata for this exact version and shown on the abstract page https://arxiv.org/abs/2606.24532v1. Full rights notice: licenses/arxiv-2606.24532-CC-BY-4.0.txt.\par\smallskip

\noindent\textit{Editorial scope.} Counted unit: the Thm whose source label is \texttt{regintN2} in the article (source0.tex lines 1089--1099, source label \texttt{regintN2}), delivered together with the article's own complete original proof of it (source0.tex lines 1101--1178, one \texttt{proof} environment of 78 lines). No mathematics is written, repaired or re-derived: statement and proof are the article's own text line for line. Required context retained uncounted: def:Lik (lines 535--537); polylog:deriv (lines 539--541); GonchL3 (lines 589--591); lem:reg\_xi1 (lines 1028--1037); lem:reg\_xi2 (lines 1057--1067). Every definition, display and label of those lines is retained unchanged. Omitted from this package and not redistributed: 1--534, 538--538, 542--588, 592--1027, 1038--1056, 1068--1088, 1100--1100, 1179--2189. Nothing inside a retained excerpt is omitted or altered: every retained body is delivered exactly as the article's lines are written, so no omission receipt of any kind accompanies this package. Citations: the retained excerpts carry no citation command at all, so no bibliography entry of the article is transcribed and no reference is redistributed. Reference adaptation: none was needed and none was made; every reference inside the delivered unit resolves inside it, so no unresolved reference or citation is produced. Printed slips kept literal: no printed word, symbol, hypothesis or number of the counted statement or of its proof was altered. Layout and class: the article's own class is \texttt{amsart} with packages including amsthm, amsmath, amssymb, amsfonts, mathrsfs, graphicx, bbm, hyperref, color, xy, comment, here, amscd, caption; the wrapper is the lane's pinned \texttt{amsart} editorial wrapper at 10pt on A4, which restates the theorem-like environments the retained text needs and adds the article's own macro definitions that the retained text needs (\texttt{\textbackslash Z}, \texttt{\textbackslash dlog}, \texttt{\textbackslash e}, \texttt{\textbackslash z}), with extra wrapper packages (bm, bbm, dsfont, xspace, cancel, mathtools). The delivered numbering is the unit's own. Assets: no retained span contains a figure environment, a graphics-inclusion command, a PDF-page inclusion, a table or a TikZ picture; no image file is copied into this package, the package declares no asset dependency and no image file is redistributed with it. Licence: version arXiv:2606.24532v1 of this article is distributed under the Creative Commons Attribution 4.0 International licence, as recorded in the arXiv licence metadata for this exact version and shown on the abstract page https://arxiv.org/abs/2606.24532v1. A full rights notice is delivered at licenses/arxiv-2606.24532-CC-BY-4.0.txt. Characters: every retained excerpt is pure ASCII, so the delivered engine, XeLaTeX with its default Latin Modern text font, reports no missing character.
\begin{equation}\label{def:Lik}
    \mathrm{Li}_k(z):=\sum_{n\geq 1} \frac{z^n}{n^k}.
\end{equation}

\begin{equation}\label{polylog:deriv}
    \mathrm{Li}_1(z)=-\log(1-z) \qquad \text{ and } \qquad \frac{d}{dz} \mathrm{Li}_k(z)=\frac{\mathrm{Li}_{k-1}(z)}{z}dz.
\end{equation}

\begin{equation}\label{GonchL3}
\mathcal{L}_3(z):=\Re\Bigl(\mathrm{Li}_3(z)-\mathrm{Li}_2(z)\log(\vert z\vert)+\frac{1}{3}\mathrm{Li}_1(z)\log(\vert z\vert)^2\Bigr)
\end{equation}

\begin{lem} \label{lem:reg_xi1}
Let $\varepsilon$ be a $2N$-th root of unity with $\varepsilon^N = -1$. Then
\begin{equation*}
    \int_{\gamma_N} r_3(2)(\xi_N^{1,1}(\varepsilon)) = \int_{\delta_N} r_3(2)(\xi)
\end{equation*}
with
\begin{equation*}
    \xi = \xi_N^{1,1}(\varepsilon) - \frac{1}{N} \sum_{r \in \Z/N\Z} \xi_N^{1,1}(\varepsilon \zeta_N^r).
\end{equation*}
\end{lem}

\begin{lem} \label{lem:reg_xi2}
    Let $\zeta$ be a root of unity. Then
    \begin{align*}
        \int_{\delta_N} r_3(2)(\xi_N^{1,1}(\zeta)) = \int_0^1 \rho_{N,\z}^1 + \rho_{N,\z}^2,
    \end{align*}
    where
    \begin{align*}
        \rho_{N,\z}^1 & = - \frac13 \log \bigl|1-\zeta (t(1-t))^{1/N}\bigr| \Bigl(-\frac{\log(1-t)}{t} - \frac{\log t}{1-t}\Bigr) dt, \\
        \rho_{N,\z}^2 & = \frac13 \log t \Bigl(-\log \bigl|1-\zeta (t(1-t))^{1/N}\bigr| \dlog(t(1-t)) + \log(t(1-t)) \dlog \bigl|1-\zeta (t(1-t))^{1/N}\bigr| \Bigr).
    \end{align*}
\end{lem}

\begin{thm}\label{regintN2}
    Let $\e$ be a $2N$-th root of unity with $\e^N=-1$. Then 
    \begin{equation*}
    \int_{\gamma_N} r_3(2)(\xi_{N}^{1,1}(\e))=\frac23  \int_0^1 \left(f_{N, \e}(t) - \frac{f_{1,-1}(t)}{N}\right) dt
      + \left(\frac1N \mathcal{L}_3(-1/4)-N^2 \mathcal{L}_3(\e (1/4)^{1/N})\right),
\end{equation*}
where $\mathcal{L}_3$ is Zagier's trilogarithm \eqref{GonchL3} and, for $x\in [0,1[$ and $z\in S^1$, the function $f_{N, z}$ is given by
\begin{equation*}
    f_{N,z}(x):=\Re(\log(1- z (x(1-x))^{1/N})) \frac{\log(1-x)}{x}.
\end{equation*}
\end{thm}

\begin{proof}
    We simplify the notation by writing $\z=\z_N$, $\xi_{N, \e}=\xi_N^{1,1}(\e)$, and $\xi_{N, \e, r}=\xi_N^{1,1}(\e \zeta_N^r)$ in this proof. By Lemma \ref{lem:reg_xi1}, we have 
\[
\int_{\gamma_N} r_3(2)(\xi_{N,\e})=\int_{\delta_N} r_3(2)(\xi_{N, \e}) - \frac{1}{N} \sum_{r \in \Z/N\Z} \int_{\delta_N} r_3(2)(\xi_{N, \e, r}).
\]
By Lemma \ref{lem:reg_xi2}, we have 
\[
\int_{\delta_N} r_3(2)(\xi_{N,\e, r})=\int_0^1 \rho_{N,\e, r}^1 + \int_0^{1} \rho_{N,\e, r}^2=\int_0^1 \rho_{N,\e, r}^1 + \int_0^{1/2} (\rho_{N,\e, r}^2 + \rho_{N,\e, r}^3),
\]
where 
\begin{align*}
    \rho_{N,\e, r}^1 & = -\frac13 \log |1- \e\zeta^r((t(1-t))^{1/N})| \Bigl( -\frac{\log (1-t)}{t} - \frac{\log t}{1-t} \Bigr) dt \\
    \rho_{N,\e, r}^2 & = \frac{1}{3} \log t \bigl( -\log |1 -\e \zeta^r((t(1-t))^{1/N}| \dlog (t(1-t)) + \log (t(1-t)) \dlog |1 -\e \zeta^r((t(1-t))^{1/N}|\bigr) \\
    \rho_{N,\e, r}^3 & = \frac{1}{3} \log (1-t) \bigl( -\log |1 -\e \zeta^r((t(1-t))^{1/N}| \dlog (t(1-t)) + \log (t(1-t)) \dlog |1 -\e \zeta^r((t(1-t))^{1/N}|\bigr).
\end{align*}
We first compute that
\begin{align*}
    \int_0^1 \rho_{N, \e}^1 & = \frac{1}{3}\int_0^1 \log|1-\e(t(1-t))^{1/N}| \left( \frac{\log(1-t)}{t} + \frac{\log(t)}{1-t} \right)dt \\ 
    & = \frac{1}{3}  \int_0^1 \Re(\log(1-\e(t(1-t))^{1/N})) \left( \frac{\log(1-t)}{t} + \frac{\log(t)}{1-t} \right)dt \\ 
    & =\frac{2}{3} \int_0^{1} \Re(\log(1 -\e(t(1-t))^{1/N})) \frac{\log(1-t)}{t} dt  \\
    & = \frac{2}{3}\int_0^{1} f_{N,\e}(t) dt.
\end{align*}
Next, we note that 
\[
\rho_{N, \e}^2+\rho_{N, \e}^3=\frac{1}{3} \log (t(1-t)) \bigl( -\log |1 -\e ((t(1-t))^{1/N}| \dlog (t(1-t)) + \log (t(1-t)) \dlog |1 -\e ((t(1-t))^{1/N}|\bigr).
\]
Putting $v=(t(1-t))^{1/N}$, we see that 
\begin{align*}
\rho_{N, \e}^2+\rho_{N, \e}^3 &=\frac{1}{3} \log (v^N) \bigl( -\log |1 -\e v| \dlog (v^N) + \log (v^N) \dlog |1 -\e v|\bigr) \\
&=\frac{1}{3} \Re\left(\log (v^N) \bigl( -\log (1 -\e v)\dlog (v^N) + \log (v^N) \dlog (1 -\e v)\bigr) \right)
\end{align*}
When $t$ ranges over $]0, 1/2[$, $v$ ranges over $]0,(1/4)^{1/N}[$, and 
we deduce that 
\begin{align*}
\int_0^{1/2} (\rho_{N, \e}^2+\rho_{N, \e}^3) & = \frac{N}{3} \Re \left(\int_{0}^{(1/4)^{1/N}} \log (v) \Bigl( -N\log (1 -\e v) \frac{dv}{v} + N\log (v) \dlog (1 - \e v)\Bigr) \right) \\
& = \frac{N^2}{3} \Re \left(\int_{0}^{(1/4)^{1/N}} \log (v) \Bigl( - \frac{\log (1 -\e v)}{v} - \frac{\e \log (v)}{1 -\e v} \Bigr) dv \right) \\
& = \frac{N^2}{3} \Re \left(\int_{0}^{(1/4)^{1/N}} g_{\e}(v) dv \right),
\end{align*}
where, for $x\in ]0,1]$ and $z\in S^1$, the function $g_{z}$ is defined by
\begin{equation}
    g_{z}(x):=- \log (x) \Bigl(  \frac{z \log (x)}{1-z x} +\frac{\log (1-z x)}{x} \Bigr).
\end{equation}
Define, for $x\in [0, 1]$ and $z\in S^1$, the function 
\begin{equation}\label{fct:Gzx}
    G_{z}(x):=-3\mathrm{Li}_3(z x)+3\mathrm{Li}_2(z x)\log(x)+\log(1-z x)(\log(x))^2,
\end{equation}
where the trilogarithm $\mathrm{Li}_3$ is defined as in \eqref{def:Lik}. One readily checks using \eqref{polylog:deriv} that $G_{z}'(x)=g_z (x)$. Moreover, since $\vert z\vert = 1$, we have
\begin{equation}\label{GtoL3}
\Re(G_z(x))=-3\mathcal{L}_3(z x),
\end{equation}
where $\mathcal{L}_3$ is Zagier's trilogarithm \eqref{GonchL3}.
It follows that 
\begin{equation*}
\int_0^{1/2} (\rho_{N, \e}^2+\rho_{N, \e}^3) = \frac{N^2}{3} \Re \left(  G_{\e}((1/4)^{1/N}) \right)=-N^2 \mathcal{L}_3(\e (1/4)^{1/N}).
\end{equation*}
Finally, we note that 
\begin{equation*}
\sum_{r\in \Z/N\Z} \log \vert 1 -\e \zeta^r (t(1-t))^{1/N} \vert = \log \left\vert \prod_{r\in \Z/N\Z} (1 -\e \zeta^r (t(1-t))^{1/N})  \right\vert  = \log \left\vert 1+t(1-t)  \right\vert.
\end{equation*}
Thus, 
\begin{align*}
\sum_{r\in \Z/N\Z} \int_0^1 \rho_{N,\e, r}^1 &= -\frac13 \int_0^1 \log (1+t(1-t)) \Bigl( -\frac{\log (1-t)}{t} - \frac{\log t}{1-t} \Bigr) dt \\ &= \frac{2}{3} \int_0^1 \log(1+t(1-t))\frac{\log(t)}{1-t} dt \\
&=\frac{2}{3} \int_0^1 f_{1,-1}(t)dt.
\end{align*}
Furthermore, we have 
\[
\sum_{r\in \Z/N\Z} (\rho_{N,\e, r}^2+\rho_{N,\e, r}^3)=\frac{1}{3} \log (t(1-t)) \bigl( -\log (1+ t(1-t)) \dlog (t(1-t)) + \log (t(1-t)) \dlog (1+ t(1-t)).
\]
Putting $u=t(1-t)$, we see that 
\[
\sum_{r\in \Z/N\Z} (\rho_{N,\e, r}^2+\rho_{N,\e, r}^3)=\frac{1}{3} \log (u) \bigl( -\log (1+ u) \dlog (u) + \log (u) \dlog (1+ u)\bigr)=\frac 13 g_{-1}(u).
\]
As $t$ ranges in $]0, 1/2[$, $u$ ranges in $]0, 1/4[$, and we obtain 
\[
\sum_{r\in \Z/N\Z} \int_0^{1/2} (\rho_{N,\e, r}^2+\rho_{N,\e, r}^3)=\frac{1}{3} \int_{0}^{1/4} g_{-1}(u) du=\frac{1}{3} G_{-1}(1/4)=-\mathcal{L}_3(-1/4).
\]
The result follows by putting everything together.
\end{proof}

\end{document}
